% \begin{abstract}
% Retrieval-augmented systems often provides a page that is true and relevant but contains only part of the answer, e.g. the question asks which disease a drug treats, and the page says \emph{hyperlipidemia}, where the answer is \emph{\textbf{familial} hyperlipidemia}. Work on sufficient context treats such pages as insufficient. We argue they form a distinct case, \emph{resolvable insufficiency}: the text cannot settle the answer, but its family name indexes a short candidate list in an ontology. Language models meet this phenomenon with \clade{} (category-level answering from deficient evidence): they answer instead of declining or looking further.  On 856 PrimeKG questions, a family-only page weakens the push toward ``not enough information'' by 13.8 points relative to a silent page (pooled over three 7--8B models, $p = 4\times10^{-5}$; significant for two of the three individually), and the freed answers are about as often wrong as right. A counterfactual edit isolates the cause: replacing the disease name with its family name lowers declines by 17.4 points relative to deleting the sentence. Confidence (AUROC 0.59--0.72 on family-only pages, against 0.73--0.87 on named pages), a 14B model (58\% precision, against 96\%) and an instruction to decline (more declines on every kind of page) all leave \clade{} in place. It replicates on species interactions from GloBI, with a taxonomy as the symbolic structure (900 questions; decline effect $-14.3$ points, $p = 10^{-9}$), where the freed answers are mostly wrong. Resolving the insufficiency is only partly solved: the ontology list holds the answer 98.5\% of the time once the family is known, but the family is found from the page only 19.8\% of the time, and a popularity prior beats both the model and literature evidence at choosing within it. The results locate where a neural reader should hand over to symbolic knowledge, and which steps of that handover remain open.
% \end{abstract}

% \begin{abstract}
% Retrieval-augmented systems often retrieve passages that are accurate and relevant, but state the answer at a coarser level of granularity than the question requires. For example, a question may ask which specific disease a drug treats (\emph{familial hyperlipidemia}), while the retrieved passage mentions the broader disease family (\emph{hyperlipidemia}). Prior work on context sufficiency treats such coarse passages simply as insufficient. We argue that they form a distinct case, which we call \emph{resolvable insufficiency}: text alone cannot settle the fine-grained answer, but the broader category it mentions maps to a short list of candidate answers in a domain ontology. When language models read these coarse, \emph{family-only} passages, they exhibit a phenomenon we call \clade{} (Category-Level Answering from Deficient Evidence): they answer instead of abstaining or searching further. Across 856 biomedical questions from PrimeKG, providing a family-only passage reduces abstention by 13.8 points compared to providing an uninformative passage. Moreover, the additional answers triggered by these coarse passages are wrong about half the time. A counterfactual experiment confirms the cause: replacing an exact disease name in a passage with its broader family name reduces abstention by 17.4 points compared to deleting the sentence. Simple fixes do not stop \clade{}. Filtering by model confidence (AUROC 0.59--0.72 on family-only passages versus 0.73--0.87 on exact-match passages), scaling to a bigger model (58\% precision versus 96\%), and explicitly instructing the model to abstain all fail to eliminate the problem. The same behavior appears in an ecological domain using species interactions from GloBI and a biological taxonomy (900 questions; abstention drop of $14.3$ points), where the extra guesses are mostly wrong. Finally, we test whether an ontology can resolve these coarse passages into exact answers. Once the correct family is known, the ontology's candidate list contains the true answer 98.5\% of the time. However, identifying the right family from the passage succeeds only 19.8\% of the time, and a simple popularity prior outperforms both the language model and literature co-mention counts when selecting the final answer. Together, our findings pinpoint where a neural reader should hand off to symbolic knowledge and which steps of that handoff remain unsolved.
% \end{abstract}

\begin{abstract}
Retrieval-augmented systems often retrieve passages that are accurate and relevant but state the answer at a coarser level of granularity than the question requires. For example, a question may ask which specific disease a drug treats (\emph{familial hyperlipidemia}), while the retrieved passage mentions the broader disease family (\emph{hyperlipidemia}). Prior work on context sufficiency treats such coarse passages simply as insufficient. We argue that they form a distinct case, which we call \emph{resolvable insufficiency}: text alone cannot settle the fine-grained answer, but the broader category it mentions maps to a short list of candidate answers in a domain ontology. When language models read these coarse, \emph{family-only} passages, they exhibit a phenomenon we call \clade{} (Category-Level Answering from Deficient Evidence): they answer instead of abstaining or searching further. On biomedical questions from PrimeKG, family-only passages reduce abstention by 13.8 points compared to uninformative passages, and the resulting guesses are wrong about half the time. A counterfactual test confirms the cause: replacing an exact disease name with its broader family name lowers abstention by 17.4 points compared to deleting the sentence. Simple fixes like confidence filtering, scaling to a larger model, and explicitly instructing the model to abstain cannot resolve \clade{}, and this replicates on ecological species interactions in GloBI. Finally, we test whether an ontology can turn coarse passages into exact answers. Given the correct family, an ontology list includes the true answer in 98.5\% of cases, but string matching finds that family on the page for only 19.8\%. Letting the model choose the family raises this to 30.0\% and end-to-end recovery from 10.0\% to 26.8\%. Overall, our results show when language models should stop guessing and consult an ontology, and highlight which steps of that process remain unsolved. Codes are available at: \url{https://anonymous.4open.science/r/CLADE}.
\end{abstract}